\chapter{Accountability}\index{Accountability}
\printchapterversion{02_Accountability}

\newthought{The Responsibility Gap}

\marginnote{%
    \begin{flushright}
    \newthought{In der Strafkolonie}\\
    Franz Kafka (1919)\\
    {\color{black}\qrcode[height=0.9cm]{https://de.wikisource.org/wiki/In_der_Strafkolonie}}
    \end{flushright}%
    \medskip
    \newthought{Minority Report}\cite{spielberg2002minority} A
    pre-crime unit arrests people before they offend, based on
    algorithmic prediction. When the system targets its own chief, the
    question of who holds the algorithm accountable becomes impossible
    to avoid.%
}

% As a leading group, you are required to moderate this session. Add margin notes in this script in case you want to add any additional questions or points to the discussion.

\newthought{A visiting traveller} witnesses a penal machine execute a condemned man for a
sentence no living person can read or explain, and is invited to endorse the system.\\

\newthought{Algorithmic governance} and black-box AI in high-stakes domains highlight
opacity and accountability gaps: audits, explainability requirements, and governance mechanisms
all attempt to close the distance between automated decision and human answerability.

\medskip
\noindent\textbf{Cases: when no one is answerable.} Kafka's machine is a
fable; the responsibility gap is not. Each system below caused real harm,
and in each the question ``who is accountable?'' had no clean answer.
\begin{itemize}
    \item \textbf{Amazon's hiring engine (2018).}\cite{dastin2018amazon}
    A model trained on a decade of engineering hires learned to downgrade
    r\'esum\'es containing the word ``women's'' and to reward male-coded
    phrasing. No one wrote the rule; the bias was inherited from the data.
    Amazon quietly scrapped the tool. Who answers for discrimination
    nobody authored?
    \item \textbf{Google Photos (2015).}\cite{simonite2018google} An image
    classifier labelled photographs of Black users as ``gorillas''. The
    remedy was not a corrected model but a deleted label: ``gorilla''
    was removed from the product. The error was hidden, not fixed.
    \item \textbf{Tesla Autopilot (2016).}\cite{ntsb2017tesla} A driver was
    killed when the system failed to distinguish a crossing trailer from
    bright sky and did not brake. Liability fractured three ways: the driver
    who trusted a feature named ``Autopilot'', the manufacturer that named
    it so, and the regulator that permitted the deployment.
    \item \textbf{The Flash Crash (2010).} Automated trading erased roughly
    a trillion dollars in minutes; no human decided to sell at scale and no
    human decided to stop. Years later a single trader was charged.
\end{itemize}

\vspace{2em}

\begin{tabular}{p{3.8cm} p{6.5cm}}
    \textbf{Required reading} &
        \textit{In der Strafkolonie}\cite{kafka1919strafkolonie} by Franz Kafka
        (1919; Reclam UB~9900) \\[3pt]
    \textbf{Preparation}      &
        Complete the Guided Reading Sheet individually before the session. \\
\end{tabular}

\vspace{3.5em}

\section*{Guided Reading Sheet}

\noindent Read \textit{In der Strafkolonie} before the session. Work through the
questions below individually, then bring your notes to the discussion.

\begin{enumerate}
    \item The condemned man never learns what he is accused of or what
    the machine will inscribe on his body. What does this tell us about
    the relationship between legibility and legitimacy? Can a system be
    just if it cannot explain itself to those it acts upon?

    \item Three figures surround the machine: the officer who
    understands and maintains it, the condemned man who cannot, and
    the visiting traveller who is invited to endorse it. How does
    responsibility shift across this triangle? Who, if anyone, is
    morally culpable for what happens?

    \item The story illustrates the idea of the \textbf{responsibility
    gap}: when a complex system causes harm, no single human actor is
    clearly at fault. Can you name a contemporary case (in law,
    medicine, finance, or infrastructure) where this gap has appeared?

    \item The EU AI Act (2024) requires human oversight for high-risk
    AI systems. What does "oversight" mean in practice if the system's
    decisions are too fast, too complex, or too numerous for humans to
    review individually?

    \item Compare \textbf{causal responsibility} (who caused the harm)
    with \textbf{moral responsibility} (who is blameworthy) and
    \textbf{legal liability} (who must compensate). Do these three always
    converge? Give an example where they do not.\\
\end{enumerate}



\section*{In-Class Experiment: The Responsibility Budget}

\noindent A short allocation exercise to anchor the debate in shared
experience. Run before the Opinion Landscape, using one concrete incident
drawn from the cases above.

\begin{itemize}
    \item \textbf{Setup.} State one incident in three sentences, for
    instance the hiring engine that filtered out a qualified applicant. On
    the board, list the actors: builder, deployer, operator, the people
    whose data trained the model, the regulator, the affected person, and
    no one.
    \item \textbf{Allocation.} Each student silently distributes one
    hundred responsibility points across those actors and writes a single
    sentence justifying their largest share. Collect the slips anonymously.
    \item \textbf{Tally.} Add the points on the board. The distribution
    will not converge: it spreads across five or six actors, and a visible
    share falls on no one. That spread is the responsibility gap, drawn from
    the room's own intuitions.
    \item \textbf{Debrief.} The two largest shares are almost always the
    provider (builder and deployer) and the user. Name them as the two
    sides of the debate to come, and have each student defend the allocation
    they made.
\end{itemize}

\section*{The EU AI Act's Risk Tiers}

\marginnote{\newthought{Risk-based regulation.} Obligations scale with the
stakes of the application, not the sophistication of the technology. A
simple model deciding creditworthiness is regulated far more tightly than
a complex one recommending films.}

\noindent The EU AI Act\cite{euaiact2024} is the first comprehensive
attempt to close the responsibility gap by statute. It sorts every system
into one of four tiers by the harm it can do, and attaches obligations that
grow with the tier. The pyramid is narrow at the top: most systems carry no
obligations.

\begin{itemize}
    \item \textbf{Unacceptable risk, prohibited.} Social scoring by public
    authorities, subliminal manipulation, untargeted scraping of faces to
    build recognition databases, and most real-time biometric
    identification in public space. These are banned outright; no amount of
    documentation or oversight makes them lawful.
    \item \textbf{High risk, permitted but heavily regulated.} Systems that
    decide access to employment, credit, education, essential services, law
    enforcement, or migration. Before deployment each requires a risk
    management system, governed training data, technical documentation,
    event logging, human oversight, and a conformity assessment audited
    against the Act.
    \item \textbf{Limited risk, transparency only.} Chatbots, emotion
    recognition, and deepfakes. The single duty is disclosure: the user
    must be told they are dealing with a machine or with synthetic content.
    \item \textbf{Minimal risk, unregulated.} Spam filters, recommendation
    engines, game AI. This is the overwhelming majority of deployed
    systems, and the Act leaves them alone.
\end{itemize}

\marginnote{\newthought{General-purpose models.} Foundation models sit on a
separate track: all owe transparency duties, and those trained above a
compute threshold carry systemic-risk obligations whatever the downstream
application.}

\newthought{Every tier is a trade.} Moving a system up the pyramid trades
speed for protection. A minimal-risk feature can be deployed at once; a
high-risk one waits on documentation, an audit, and a human-oversight
design: months of work before deployment. The Act judges that where a
system decides who gets a job, a loan, or bail, the price is worth paying.
Logging lets a named party be held to account; human oversight keeps a
person, not a dataset, the owner of the decision; the conformity assessment
proves rigor before harm rather than after. Developers object: the same
rigor slows European deployment against faster
jurisdictions, and mandated transparency exposes model internals that are
also security and competitive assets. Speed, user rights, security, and
development rigor cannot all be maximised at once. The tiers fix the
exchange rate between them, and the debate below turns on whether the Act
has set that rate correctly.

\section*{Opinion Landscape and Debate}

\noindent \textbf{Format:} Surrounded-style debate. Follow the shared
debate protocol for the speaker's list rules and moderator scripts.

\medskip
\noindent \textbf{Opening question:} When an AI system causes harm, who
takes responsibility: the builder and deployer who put it into the world,
or the user who chose to apply it to a purpose?

\marginnote{\newthought{Instrument or agent?} The law treats the two
differently: you answer for a tool you wield; a maker answers for what it
sends into the world. An AI sits between the two, and the debate turns on
which side it falls.}

\newthought{Opinion~A, User Accountability.}\quad Responsibility rests
with the user who put the system to use. Builder and deployer supply a
capability; the user aims it at a purpose, and that choice carries the
responsibility.
\vspace{1.5em}


\medskip
\noindent\textbf{The Builder} insists the system performs within documented
    parameters and that responsibility follows its use.
    \begin{itemize}
        \item \textit{Our documentation states that human review is
        mandatory before any recommendation is acted upon. The user relied
        on the output without it. We cannot be held liable for how our tool
        was used.}
        \item \textit{A tool has no purpose until someone picks it up and
        gives it one. Using a tool for a purpose is an action, and the
        person who chose the purpose owns the outcome. We do not blame the
        scalpel for the surgeon's cut.}
        \item \textit{Extend our liability to every downstream use and no
        one will build decision-support tools at all.}
    \end{itemize}

\medskip
\noindent\textbf{The Deployer} fielded the system and argues the final
    call sat with the user.
    \begin{itemize}
        \item \textit{We procured a system certified for exactly this task,
        and it performed within specification. Our role is to make the
        capability available, not to make each individual decision with it.}
        \item \textit{The final decision sat with the user. They read the
        output and chose to act on it. That choice, not our procurement, is
        the proximate cause of the harm.}
        \item \textit{The user is a trained professional exercising
        discretion, and discretion carries responsibility. We cannot stand
        behind every judgement each operator makes.}
    \end{itemize}



\newthought{Opinion~B, Provider Accountability.}\quad Responsibility rests with those who
built and fielded the system. An AI is not a neutral instrument the user
fully controls; it is closer to something its makers raised, and the user
cannot answer for a decision they were never in a position to understand.
\vspace{1.5em}


\medskip
\noindent\textbf{The User} refuses to answer for a black box they could
    neither open nor overrule.
    \begin{itemize}
        \item \textit{You handed me a black box and told me to trust it. I
        cannot open it, cannot audit the weights, cannot see why it decided
        what it decided. You cannot make a decision unauditable and then
        blame me for it.}
        \item \textit{You call it a tool I chose to use, but a hammer does
        not learn. You trained this on data you chose and tuned it toward
        behaviour you selected. That is closer to raising a child than
        forging an instrument, and a parent answers for the child they
        shaped.}
        \item \textit{The knowledge and the control both sat with the people
        who built and fielded it. Responsibility should sit where they do,
        not with the last human in the chain.}
    \end{itemize}

\medskip
\noindent\textbf{The Regulator} argues that whoever profits from and
    controls the design must carry the risk.
    \begin{itemize}
        \item \textit{Whoever captures the upside must internalise the
        downside. A disclaimer buried in a manual cannot transfer a hazard
        the user has no means to evaluate.}
        \item \textit{Product-safety law already binds those who build and
        field a product to its foreseeable real-world use, not to the
        idealised use their documentation describes.}
        \item \textit{Placing the whole burden on the last human to touch
        the system builds a moral crumple zone\cite{elish2019crumple}: a
        person positioned to absorb the blame while powerless to prevent the
        harm. That is not accountability; it is its appearance.}
    \end{itemize}


\newthought{Key Learnings}
The responsibility gap, the structural absence of a single culpable
actor in distributed sociotechnical systems, is not caused by
negligence but by the architecture of those systems: standard
categories of fault (criminal intent, professional negligence, product
liability) apply only awkwardly when causation is diffuse. Closing it
requires deliberate governance design, not post-hoc blame. In practice the
gap becomes a contest over allocation: builder and deployer point to
documented limits and the user's discretion; the user points to an opaque
product they could neither audit nor overrule. Beneath it lies one
unsettled question: is an AI an instrument, whose user answers for its use,
or a trained agent, whose makers answer for its formation? Explainability, the
property of being able to account for outputs in terms a human
decision-maker can evaluate, is a necessary but not sufficient
condition for accountability; a system can be interpretable and still
have no human willing to own its outputs. Meaningful human control, a
regulatory concept that requires humans to retain the ability to
understand, intervene in, and reverse AI decisions, is contested in
practice: it may mean the ability to intervene, to understand, to
reverse, or merely to authorise, and these are different requirements
with different institutional costs. The lessons of aviation safety culture, medical liability law,
and financial auditing all offer partial models for algorithmic governance, but none transfers cleanly.

\medskip
\noindent\textit{Teaser: what happens to accountability when the system
decides faster than anyone can follow or override?}

\clearpage
